Sur ces images, l'intensité lumineuse obtenue est exprimée en LSB (Lower Significative Bit, unité de mesure spécifique de la caméra). On indique qu'un éclair traditionnel provoque un niveau de 836 LSB sur la caméra opérant dans le visible et donne 30 LSB sur la caméra filtrée.

\TLEFigure{4}{Enregistrements effectués depuis l'ISS par la microcaméra filtrée à gauche et par la microcaméra dans le visible à droite. Un trait noir isole la détection d'un sylphe sur l'image de droite - Figures toujours tirées de l'article référencé sur la figure 3.}

\begin{itemize}
    \item[3.] Montrer que l'analyse des deux images de la figure 4 prouve bien la détection d'un sylphe.
    \item[4.] L'ISS se situe à une altitude $h_{\text{iss}}$ d'environ 400 km. Montrer en utilisant le théorème de GAUSS, et moyennant certaines hypothèses que l'on précisera, que le champ gravitationnel qu'elle subit est de la forme :
    \begin{equation*}
        G_{\text{iss}} = g_0 \left( 1 + \frac{h_{\text{iss}}}{R_T} \right)^{-2}
    \end{equation*}
    où $g_0 = 9,8 \text{ m} \cdot \text{s}^{-2}$ est le champ de pesanteur à la surface de la Terre.
    \item[5.] Après avoir montré qu'elle est constante, déterminer l'expression $v_{\text{iss}}$ de la norme de la vitesse de l'ISS sur son orbite circulaire. Calculer sa valeur numérique en $\text{km} \cdot \text{s}^{-1}$.
    \item[6.] Déterminer l'expression de la période de révolution $T_{\text{iss}}$ de l'ISS autour de la Terre et faire le lien avec la troisième loi de KÉPLER.
    \item[7.] Les deux microcaméras de l'ISS ont pour objectif d'enregistrer, lors d'un passage au-dessus d'un orage, la totalité de l'évolution d'un sylphe. Commenter la possibilité et la précision de ces enregistrements. On justifiera les réponses de façon numérique.
\end{itemize}

Il est difficile d'observer des sylphes depuis la Terre dans le visible car le dioxygène $\text{O}_2$ présent dans l'atmosphère absorbe fortement à 762 nm. Toutefois, le dioxygène se raréfie fortement avec l'altitude, ce qui permet des observations au-delà de 40 km d'altitude comme cela est le cas avec l'ISS. Nous allons étudier l'évolution de la pression partielle $P_{\text{O}_2}(z)$ de dioxygène en fonction de l'altitude $z$. Le modèle que nous utilisons est le suivant :

\begin{itemize}
    \item L'atmosphère est assimilée à un gaz parfait de masse molaire $M = 29 \text{ g} \cdot \text{mol}^{-1}$.
    \item La fraction molaire de dioxygène $\text{O}_2$ est supposée la même à toute altitude $z$ et égale à celle à la surface de la Terre à savoir $x_{\text{O}_2} = 20\%$.
\end{itemize}
